\section{Beyond Canonical Ensemble}

\subsection{Internal Energy Minimization}

\paragraph{Two ways to slice the entropy landscape.}
Consider the entropy $S(U,X)$ as a surface over the $(X,U)$ plane, where $X$
is an extensive variable that can be freely partitioned, such as a volume
shared between two bodies. Draw the \textcolor{blue}{level curves} of $S$,
the curves along which $S$ is constant, like contour lines on a topographic
map. The map can be sliced in two ways:
\begin{itemize}
    \item Entropy maximum at fixed $U$: walk along a horizontal line,
    $U=\text{const}$, and find the highest contour it reaches. At that point
    the line touches the contour without crossing it, so the contour is
    tangent to the line.
    \item Energy minimum at fixed $S$: walk along one contour,
    $S=\text{const}$, and find its lowest point. There the contour has a
    horizontal tangent.
\end{itemize}
Both slices therefore select a point where the level curve of $S$ has a
horizontal tangent in the $(X,U)$ plane.

\paragraph{One geometric condition, two descriptions.}
The gradient
\begin{equation*}
    \nabla S = \left(\frac{\partial S}{\partial X},\
    \frac{\partial S}{\partial U}\right)
\end{equation*}
is perpendicular to the level curve. If $\partial S/\partial X=0$ at a point,
the gradient points purely along the $U$-direction, so the level curve points
purely along the $X$-direction: it is horizontal. A horizontal tangent means
$dU/dX|_S=0$, a stationary point of $U$ along the contour. The implicit
function theorem makes this explicit:
\begin{equation}
    \boxed{\frac{\partial U}{\partial X}\bigg|_S
    = -\,\frac{\partial S/\partial X\,|_U}{\partial S/\partial U\,|_X}.}
\end{equation}
Since $\partial S/\partial U=1/T$ is finite and nonzero, the left side
vanishes exactly when the numerator vanishes:
\begin{equation}
    \boxed{\frac{\partial S}{\partial X}\bigg|_U=0
    \quad\Longleftrightarrow\quad
    \frac{\partial U}{\partial X}\bigg|_S=0.}
\end{equation}
The two conditions are not coincidentally equal. They describe the same
point on the same contour line.

\paragraph{Generality of the argument and the role of stability.}
This argument does not depend on entropy or energy. For any smooth function
$f(x,y)$ that is monotonic in $y$ at fixed $x$, the points where $f$ is
extremal in $x$ at fixed $y$ coincide with the points where $y$ is extremal
in $x$ along a level set of $f$, because both say that the contour is flat.
Thermodynamics is the case $f=S$ and $y=U$, with $x$ any extensive variable;
monotonicity follows from $\partial S/\partial U=1/T>0$. What the
thermodynamic stability condition, $\partial^2S/\partial X^2<0$, adds is the
type of flat point: a maximum of $S$ at fixed $U$ corresponds to a minimum of
$U$ at fixed $S$.

\paragraph{Which flat point: maximum of $S$, minimum of $U$.}
We differentiate the implicit-function relation once more along a contour of
constant $S$. Using $\partial U/\partial X|_S=-T\,\partial S/\partial X|_U$,
\begin{equation*}
    \frac{\partial^2U}{\partial X^2}\bigg|_S
    =-T\,\frac{d}{dX}\bigg|_S\left(\frac{\partial S}{\partial X}\bigg|_U\right)
    -\frac{\partial T}{\partial X}\bigg|_S\frac{\partial S}{\partial X}\bigg|_U.
\end{equation*}
At the stationary point, $\partial S/\partial X|_U=0$ and
$\partial U/\partial X|_S=0$. The second term then vanishes, and in the first
term the chain rule reduces to $\partial^2S/\partial X^2|_U$:
\begin{equation}
    \boxed{\frac{\partial^2U}{\partial X^2}\bigg|_S
    =-T\,\frac{\partial^2S}{\partial X^2}\bigg|_U.}
\end{equation}
Since $T>0$ and $\partial^2S/\partial X^2<0$ by stability, the right side is
positive. A maximum of $S$ at fixed $U$ is therefore a minimum of $U$ at
fixed $S$.

\subsection{Thermodynamic Potential}

\paragraph{Entropy as the microcanonical thermodynamic potential.}
In the microcanonical ensemble, the entropy $S(U,V,N)$ is a function of
state: the values of $U$, $V$, and $N$ determine $S$ completely.
Differentiating $S$ yields every conjugate intensive variable, and $S$ is the
extremized quantity of the ensemble. For these reasons, $S$ is called the
\textcolor{blue}{thermodynamic potential} of the ensemble, and $(U,V,N)$ are
its \textcolor{blue}{natural variables}. A potential $\Psi(x_1,\ldots,x_n)$
has natural variables $(x_1,\ldots,x_n)$ when:
\begin{enumerate}
    \item[(i)] differentiating $\Psi$ with respect to each $x_i$ recovers
    every conjugate variable: $d\Psi$ is an exact differential in exactly
    these variables, with nothing missing and nothing required from outside;
    \item[(ii)] the $x_i$ are precisely the quantities externally fixed or
    controlled under the physical condition that $\Psi$ describes, that is,
    the hard constraints of the ensemble, not an arbitrary choice of
    variables.
\end{enumerate}

\paragraph{The need for a potential with natural variables $(N,V,T)$.}
Because $S$ is strictly increasing in $U$, it can be inverted to give
$U(S,V,N)$, which is also a function of state. Its derivatives likewise
yield $T$, $P$, and $\mu$, and $U$ is minimized at fixed $(S,V,N)$. In the
canonical ensemble, however, $U$ is no longer a controlled parameter; the
temperature is. The entropy is then known only as a consequence of $T$, not
as a coordinate we can vary at will. We therefore need a potential that
preserves the full information of the system and uses the natural variables
of the ensemble, $(N,V,T)$.

\subsection{The Legendre Transformation}

\paragraph{Constructing the transform from tangent lines.}
\textcolor{red}{In this section we denote the slope by $\xi$ instead of $p$,
to avoid confusion with the probability $p_i$ and the pressure $P$.}
Let $f(x)$ be a function. At each point $x$, the tangent line has slope
$\xi(x)=f'(x)$ and intercept $f(x)-\xi(x)\,x$. Expressing the intercept as a
function of the slope defines the \textcolor{blue}{Legendre transform}
\begin{equation}
    \boxed{g(\xi)=f(x)-\xi\,x,\qquad \xi=f'(x).}
\end{equation}
We say that $g(\xi)$ conserves the information of $f(x)$ if $f$ can be
reconstructed from $g$ alone.

\paragraph{Reconstructing $f$ from its transform.}
Each tangent line, labeled by its slope $\xi$, is $y=\xi x+g(\xi)$, so the
tangent lines form a family with one line per $\xi$. Fix a point $x$ and
consider the vertical gap between this family and the curve,
\begin{equation*}
    \Delta(\xi,x)=\xi x+g(\xi)-f(x).
\end{equation*}
A function $f$ is \textcolor{blue}{convex} if $-f$ is concave, that is, if
the inequality in the definition of concavity is reversed. If $f$ is convex,
every tangent line lies on or below the curve, so $\Delta\leq0$ for all
$\xi$, with equality exactly for the line tangent at $x$. If $f$ is concave,
every tangent line lies on or above the curve, so $\Delta\geq0$ with the same
equality condition. In either case, the tangent line is a genuine extremum of
$\Delta$ over $\xi$, not merely a stationary point:
\begin{equation*}
    \frac{\partial\Delta}{\partial\xi}=x+g'(\xi)=0
    \quad\Longleftrightarrow\quad x=-g'(\xi).
\end{equation*}
Calling the solution $\xi^*(x)$, we reconstruct $f$ as
\begin{equation}
    \boxed{x=-g'(\xi),\qquad f(x)=\xi^*(x)\,x+g\big[\xi^*(x)\big].}
\end{equation}
Hence $f$ and $g$ determine each other, and no information is lost. The
transformation trades the variable $x$ for its conjugate $\xi=f'(x)$.

\subsection{Multivariable Legendre Transformation}

\paragraph{Extending the transform to several variables.}
Let $f(x_1,\ldots,x_n)$ be a function of $\mathbf{x}=(x_1,\ldots,x_n)$. At
each point, the \textcolor{blue}{tangent hyperplane} has gradient
$\boldsymbol{\xi}(\mathbf{x})=\nabla f(\mathbf{x})$, with components
$\xi_i=\partial f/\partial x_i$, and intercept $f-\boldsymbol{\xi}\cdot\mathbf{x}$.
The Legendre transform is
\begin{equation}
    \boxed{g(\boldsymbol{\xi})
    = f(\mathbf{x})-\boldsymbol{\xi}\cdot\mathbf{x}
    = f(\mathbf{x})-\sum_{i=1}^n\xi_i x_i.}
\end{equation}

\paragraph{Reconstructing $f$ from its transform.}
Each tangent hyperplane, labeled by its gradient $\boldsymbol{\xi}$, is
$y=\boldsymbol{\xi}\cdot\mathbf{x}+g(\boldsymbol{\xi})$, so the hyperplanes
form a family with one member for each $\boldsymbol{\xi}\in\mathbb{R}^n$.
Fix a point $\mathbf{x}$ and consider the vertical gap between this family
and the surface,
\begin{equation*}
    \Delta(\boldsymbol{\xi},\mathbf{x})
    =\boldsymbol{\xi}\cdot\mathbf{x}+g(\boldsymbol{\xi})-f(\mathbf{x}).
\end{equation*}
If $f$ is convex, every tangent hyperplane lies on or below the surface, so
$\Delta\leq0$ for all $\boldsymbol{\xi}$, with equality exactly for the
hyperplane tangent at $\mathbf{x}$. If $f$ is concave, $\Delta\geq0$ with the
same equality condition. In either case, that is, when the Hessian of $f$ is
semidefinite everywhere, the tangent hyperplane is a genuine extremum of
$\Delta$ over $\boldsymbol{\xi}$:
\begin{equation*}
    \frac{\partial\Delta}{\partial\xi_i}
    =x_i+\frac{\partial g}{\partial\xi_i}=0
    \quad\Longleftrightarrow\quad
    x_i=-\frac{\partial g}{\partial\xi_i}
    \qquad(i=1,\ldots,n).
\end{equation*}
Calling the solution $\boldsymbol{\xi}^*(\mathbf{x})$, we reconstruct $f$ as
\begin{equation}
    \boxed{x_i=-\frac{\partial g}{\partial\xi_i},\qquad
    f(\mathbf{x})=\boldsymbol{\xi}^*(\mathbf{x})\cdot\mathbf{x}
    +g\big[\boldsymbol{\xi}^*(\mathbf{x})\big].}
\end{equation}
No information is lost: the transformation trades the variables
$x_1,\ldots,x_n$ for their conjugates $\xi_i=\partial f/\partial x_i$.

\paragraph{Partial Legendre transforms.}
In thermodynamics, one often transforms only a subset
$I\subseteq\{1,\ldots,n\}$ of the variables and leaves the rest untouched:
\begin{equation}
    \boxed{g\big(\{\xi_i\}_{i\in I},\{x_j\}_{j\notin I}\big)
    =f(\mathbf{x})-\sum_{i\in I}\xi_i x_i,
    \qquad \xi_i=\frac{\partial f}{\partial x_i}\bigg|_{\mathbf{x}}.}
\end{equation}
The construction above goes through one variable at a time. Each $x_i$ with
$i\in I$ is swapped for its conjugate $\xi_i$, and each $x_j$ with
$j\notin I$ is held fixed as a spectator.

\paragraph{Convexity of partial transforms.}
If $f$ is convex, a partial Legendre transform is concave in the transformed
variables $\xi_i$ and convex in the untransformed variables $x_j$. Since
$U(S,V,N)$ is convex by thermodynamic stability, $F(T,V,N)$ is concave in
$T$ and convex in $V$, and $H(S,P,N)$ is concave in $P$ and convex in $S$.
This is why $F$ is minimized over $V$ at fixed $T$, but not over $T$.

\subsection{$NVT$ Ensemble - Helmholtz Free Energy}

\paragraph{Helmholtz free energy as the Legendre transform of $U$.}
The thermodynamic potential of the \textcolor{blue}{$NVT$ ensemble}, or
canonical ensemble, is the \textcolor{blue}{Helmholtz free energy}. It is the
partial Legendre transform of $U(S,V,N)$ that trades $S$ for its conjugate
$T=\partial U/\partial S|_{V,N}$:
\begin{equation*}
    F(N,V,T)=U-\frac{\partial U}{\partial S}\bigg|_{V,N}S,
\end{equation*}
that is,
\begin{equation}
    \boxed{F(N,V,T)=U-TS.}
\end{equation}

\paragraph{Conjugate variables of $F$.}
We check that $F$ generates every conjugate variable. Since $U$ depends on
$S$, and $S$ in turn depends on $T$, $V$, and $N$, each derivative of $F$
must include the contribution of $S$:
\begin{align*}
    \frac{\partial F}{\partial T}\bigg|_{V,N}
    &=\frac{\partial U}{\partial S}\frac{\partial S}{\partial T}-S
    -T\frac{\partial S}{\partial T}=-S,\\
    \frac{\partial F}{\partial V}\bigg|_{T,N}
    &=\frac{\partial U}{\partial V}+\frac{\partial U}{\partial S}\frac{\partial S}{\partial V}
    -T\frac{\partial S}{\partial V}
    =-P+T\frac{\partial S}{\partial V}-T\frac{\partial S}{\partial V}=-P,\\
    \frac{\partial F}{\partial N}\bigg|_{T,V}
    &=\frac{\partial U}{\partial N}+\frac{\partial U}{\partial S}\frac{\partial S}{\partial N}
    -T\frac{\partial S}{\partial N}
    =\mu+T\frac{\partial S}{\partial N}-T\frac{\partial S}{\partial N}=\mu.
\end{align*}
Hence
\begin{equation}
    \boxed{\frac{\partial F}{\partial T}\bigg|_{V,N}=-S,\qquad
    \frac{\partial F}{\partial V}\bigg|_{T,N}=-P,\qquad
    \frac{\partial F}{\partial N}\bigg|_{T,V}=\mu.}
\end{equation}
$F$ behaves much like $U$: it has the same dimension and yields the same $P$
and $\mu$, while $T$ has become an argument and $-S$ its conjugate.

\paragraph{Minimization of $F$ from entropy maximization.}
The same result follows from MaxEnt applied to the combined system. Let the
system, at fixed $V$ and $N$, exchange energy with a reservoir at temperature
$T$, and treat the pair as isolated, as in the zeroth-law setup. The
reservoir gains the energy that the system loses, so
$dS_{\text{res}}=-dU/T$, and the total entropy changes by
\begin{equation}
    \boxed{dS_{\text{tot}}=dS-\frac{dU}{T}=-\frac{dF}{T}\geq0.}
\end{equation}
Maximizing the total entropy is therefore equivalent to minimizing $F$ of the
system alone. The factor $-TS$ in $F$ accounts for the reservoir.

\paragraph{Free energy and the partition function.}
Substituting $S=k_B(\beta U+\ln Z)$ and $\beta=1/(k_BT)$ into $F=U-TS$ gives
\begin{align*}
    F=U-\frac{1}{k_B\beta}\,k_B\left(\beta U+\ln Z\right)
    =U-U-\frac{\ln Z}{\beta},
\end{align*}
so
\begin{equation}
    \boxed{F=-k_BT\ln Z.}
\end{equation}

\subsection{$NPS$ Ensemble - Enthalpy}

\paragraph{Enthalpy as the Legendre transform of $U$ in the volume.}
For some systems, it is more natural to control the pressure than the
volume, because pressure is intensive and does not depend on system size. We label this the \textcolor{blue}{$NPS$ ensemble}, after the natural
variables of $H$. Its thermodynamic potential
is the \textcolor{blue}{enthalpy}, the partial Legendre transform of
$U(S,V,N)$ that trades $V$ for its conjugate $-P=\partial U/\partial V|_{S,N}$:
\begin{equation*}
    H(N,P,S)=U-\frac{\partial U}{\partial V}\bigg|_{S,N}V,
\end{equation*}
that is,
\begin{equation}
    \boxed{H(N,P,S)=U+PV.}
\end{equation}
It generates every conjugate variable:
\begin{equation}
    \boxed{\frac{\partial H}{\partial S}\bigg|_{P,N}=T,\qquad
    \frac{\partial H}{\partial P}\bigg|_{S,N}=V,\qquad
    \frac{\partial H}{\partial N}\bigg|_{S,P}=\mu.}
\end{equation}

\paragraph{Enthalpy as exchanged heat, and its minimization.}
By the first law, at constant $(N,P)$,
\begin{equation}
    \boxed{\delta Q=dU+P\,dV=dH.}
\end{equation}
The Clausius inequality, $\delta Q\leq T\,dS$, then gives $dH\leq T\,dS$, and
at constant entropy
\begin{equation}
    \boxed{dH\leq0.}
\end{equation}
Any spontaneous process at fixed $(N,P,S)$ therefore proceeds toward lower
$H$ and halts when $H$ reaches its minimum, which recovers the equilibrium
condition established earlier. Most processes are not isentropic, however,
so enthalpy is mainly used as the heat exchanged at fixed pressure. For
example, it measures the heat of a phase transition in chemistry.

\paragraph{Labels and conventional names of ensembles.}
\textcolor{red}{From here on, we label each ensemble by the natural variables
of its thermodynamic potential, such as $NVT$ for $F$ and $NPS$ for $H$.
The conventional name is given alongside.}
The conventional label of an ensemble lists the quantities held fixed as
constraints. For $NPS$, the conventional name is the
\textcolor{blue}{isoenthalpic-isobaric ensemble}, $NPH$, which fixes $N$,
$P$, and the enthalpy $H$. It relates to $H(N,P,S)$ as the microcanonical
ensemble $NVU$ relates to $U(S,V,N)$: the same state, described either by the
fixed energy-like quantity or by the entropy. The ensembles in this chapter
and the next are:
\begin{center}
\begin{tabular}{lll}
    Natural variables & Potential & Conventional name \\ \hline
    $(S,V,N)$ & $U$ & microcanonical ($NVU$) \\
    $(T,V,N)$ & $F$ & canonical ($NVT$) \\
    $(S,P,N)$ & $H$ & isoenthalpic-isobaric ($NPH$) \\
    $(T,V,\mu)$ & grand potential & grand canonical ($\mu VT$) \\
    $(T,P,N)$ & $G$ & isothermal-isobaric ($NPT$) \\
\end{tabular}
\end{center}